% Chapter 9: original source Spin/SU2 real 48-label commutation and physical domain gap.
% Frozen green: C48-0 #42779878513; C48-0B #42787913738; C48-1A0 #42798054277.
\chapter[Gauge Compatibility and Physical Sector Reduction]{Gauge Compatibility, Sixteen-Charge Covariance, and the Physical Reducing-Sector Problem}
\label{ch:covariance}

\section{Two symmetries sharing the same Majorana labels}
The original matrix system has a spatial symmetry acting on spinor
indices and a color gauge symmetry acting on three adjoint labels.
Their linear actions on real Majorana labels are expected to commute,
but a mathematically correct proof must identify both with the
\emph{same original} 48-label ordering, including the original
$\mathrm{SU}(2)$ color coefficients. This is more subtle than proving
that two formally written Kronecker products commute: the two source
constructions originally came through different combinations of the
24 complex creation/annihilation modes and their 48 real parts.

Let $q\in\mathrm{SU}(2)$ be a source gauge element. Write
$R_3(q)$ for its genuine adjoint-color matrix on
$\mathbb R^3$. Let $D_{16}(g)$ be the source spinor action.
The original source reindexing $E$ from Chapter~\ref{ch:spin-geometry}
provides the two real forty-eight-label actions
\begin{equation}\label{eq:spatial-vs-color-labels}
 S_g=E\bigl(I_3\otimes D_{16}(g)\bigr)E^{-1},\qquad
 C_q=E\bigl(R_3(q)\otimes I_{16}\bigr)E^{-1}.
\end{equation}
These are matrices on $V_{48}$; neither is a Hilbert-space operator
on the full bosonic--fermionic $L^2$ state space.

\begin{theorem}[Original 48-Majorana label commutation; strict-qualified]
\label{thm:c48-commutation}
For all original source $g\in G_9$ and $q\in\mathrm{SU}(2)$,
\begin{equation}\label{eq:label-commutation}
 S_gC_q=C_qS_g.
\end{equation}
The source $C_q$ also agrees entry-for-entry with the independent
original S1J gauge Majorana matrix obtained from the 24-mode Fock
construction; hence \eqref{eq:label-commutation} refers to the
actual original gauge labels.
\end{theorem}
\begin{proof}[Two-step source proof]
First calculate in the original paired color--spin tensor ordering:
\[
 (I_3\otimes D_{16})(R_3\otimes I_{16})
   =(I_3R_3)\otimes(D_{16}I_{16})=R_3\otimes D_{16},
\]
and the reverse product is the same. Since $E$ is an exact source
reindexing ring equivalence, its transport preserves products.
This proves the C48-0 matrix commutator on all 48 real labels.

The C48-0B proof separately expands the original color/real-imaginary
source indices and compares the resulting entries with the S1J
matrix. It invokes the accepted G3B one-particle coefficient theorem,
so the identification is an exact equality of two independently
constructed \emph{original} source matrices, not a definitional
renaming. Consequently the label commutator also holds against the
existing genuine S1J source gauge matrix.
\end{proof}
\formalrecord{\path{C48C0OriginalSpinGaugeMajoranaCommutation.lean},
\path{source9SpinOriginalPaired48_commutes_sourceGauge};
strict-green \href{https://github.com/etblink/math/actions/runs/42779878513}{\#42779878513}.
\path{C48C0BOriginalSourceGaugeS1JMatrixIdentity.lean},
\path{sourceGaugeOriginalPaired48_eq_sourceRealMajoranaMatrix};
strict-green \href{https://github.com/etblink/math/actions/runs/42787913738}{\#42787913738}.
These receipts supersede earlier notes that called the gates pending.}

\section{From the label commutator to active Fock-generator compatibility}
A matrix equality is not yet an operator equality, but it provides
exactly the algebra needed to compare two operations on the original
active Majorana map $\Gamma:V_{48}\to\mathscr A$.
For the source gauge matrix $C_q$ define
$\Gamma_q(v):=\Gamma(C_qv)$. The previously qualified
\eqref{eq:active-covariance} immediately suggests
\begin{equation}\label{eq:one-sided-gamma-compatibility}
 \alpha_g(\Gamma(C_qv))
 =\Gamma(S_gC_qv)=\Gamma(C_qS_gv),
 \qquad v\in V_{48}.
\end{equation}
The actual C48-1A0 proof constructs this one-sided identity with the
source S1J matrix and independently checks it on the original labels.

\begin{theorem}[Original active-gamma/S1J compatibility; strict-qualified]
\label{thm:c48-active-compat}
The source Spin(9) Fock-operator automorphism commutes with the
\emph{source S1J gauge transformation on every original active
normalized Majorana generator}, in the explicit one-sided sense
\eqref{eq:one-sided-gamma-compatibility}.
\end{theorem}
\begin{proof}[Source dependency explanation]
Apply Theorem~\ref{thm:active-majorana} to the real vector $C_qv$.
Use Theorem~\ref{thm:c48-commutation} to interchange its spin and
gauge matrices. The result is exactly the active Majorana map of the
source gauge transform of $S_gv$. The Lean proof verifies the concrete
S1J coefficient and real-Majorana carrier through the C48-0B bridge,
rather than introducing an anonymous commuting pair of matrices.
\end{proof}
\formalrecord{\path{C48C1A0OriginalActiveGammaS1JSpinCompatibility.lean},
\texttt{spin9\_activeGamma\_S1J\_gauge\_commute}; strict-green
\href{https://github.com/etblink/math/actions/runs/42798054277}{\#42798054277}.
This accepted scope is on original gamma operators.}

\section{Why the genuine gauge unitaries need another theorem}
The true original gauge symmetry on the Fock-vector fiber is carried
by the source \texttt{fermionGaugeUnitaryHom} and the
\texttt{GaugeData} of Chapter~\ref{ch:gauge}. In the operator algebra,
its action is conjugation by $U_q$. We ultimately need the spatial
$\alpha_g$ to preserve this genuine gauge representation, or at least
to preserve its physical fixed projection.

The distinction can be written precisely. Even if one has proved
\begin{equation}\label{eq:automorphism-commute-condition}
 \alpha_g\circ\operatorname{Ad}_{U_q}
   =\operatorname{Ad}_{U_q}\circ\alpha_g
 \quad\text{on }\mathscr A,
\end{equation}
this does not immediately imply $\alpha_g(U_q)=U_q$. Indeed,
\eqref{eq:automorphism-commute-condition} would imply
\[
 \operatorname{Ad}_{\alpha_g(U_q)}
       =\operatorname{Ad}_{U_q},
\]
so the ratio $\alpha_g(U_q)U_q^{-1}$ commutes with every element of
the full matrix algebra. Its center is $\mathbb C I$, hence
\begin{equation}\label{eq:phase-ambiguity}
 \alpha_g(U_q)=c_g(q)U_q,\qquad |c_g(q)|=1.
\end{equation}
There is a useful refinement: \emph{once the full operator-level
commutation hypothesis is established for the genuine source
unitary homomorphism, the phase cannot remain nontrivial}. It is
automatically a one-dimensional character of the actual
$\mathrm{SU}(2)$ group, which has no nontrivial characters.
The missing source theorem is therefore the full automorphism
compatibility and its precise binding to the original unitary
carrier, rather than a freely adjustable independent scalar phase.

\begin{proposition}[Gauge-character cancellation; conditional
conventional argument]\label{prop:su2-phase-cancellation}
Let $U:\mathrm{SU}(2)\to U(\Fermion)$ be a genuine unitary
group homomorphism and let $\alpha$ be a complex star-algebra
automorphism of $\mathscr A$ satisfying
\[
 \alpha\circ\operatorname{Ad}_{U_q}
 =\operatorname{Ad}_{U_q}\circ\alpha
 \qquad\text{for every }q\in\mathrm{SU}(2).
\]
Then $\alpha(U_q)=U_q$ for every $q\in\mathrm{SU}(2)$.
\end{proposition}
\begin{proof}
The displayed commutation gives
$\operatorname{Ad}_{\alpha(U_q)}
=\operatorname{Ad}_{U_q}$. Since the center of the full
finite matrix algebra is $\mathbb C I$, and $\alpha$
preserves the star, there is $c(q)\in U(1)$ with
$\alpha(U_q)=c(q)U_q$. Because $U$ is a homomorphism and
$\alpha$ preserves multiplication,
\[
 c(qr)U_{qr}=\alpha(U_qU_r)
 =c(q)c(r)U_qU_r,
\]
hence $c(qr)=c(q)c(r)$. Thus $c$ is a group character.
The group $\mathrm{SU}(2)$ is perfect: every element is
conjugate to a diagonal matrix
$\operatorname{diag}(e^{it},e^{-it})$, and each such matrix
is a commutator. For example, with
$a=\operatorname{diag}(e^{it/2},e^{-it/2})$ and
$j=\left(\begin{smallmatrix}0&1\\-1&0\end{smallmatrix}\right)$,
one has $aja^{-1}j^{-1}=a^2$. Every character into
the abelian group $U(1)$ is therefore trivial.
So $c(q)=1$.
\end{proof}
This is a \emph{conditional} conventional proposition; it does
not supply the unproved full source
\eqref{eq:automorphism-commute-condition} nor an accepted
Lean gauge-unitary fixedness theorem. It uses the genuine
$\mathrm{SU}(2)$ homomorphism law, not arbitrary phases
attached independently to each gauge element.

\begin{openproblem}[Gauge-unitary fixedness and physical projection]
Establish the full source operator-level compatibility
\eqref{eq:automorphism-commute-condition} on the \emph{actual}
Fock operator algebra with the \emph{actual} source
$\mathrm{SU}(2)$ gauge unitary homomorphism. Then apply
Proposition~\ref{prop:su2-phase-cancellation}, prove the fixedness
of the genuine original fermionic gauge projection, and bind those
facts to the full-space combined boson--Fock action. The conditional
character argument alone is not the missing source compatibility theorem.
\end{openproblem}

\section{The missing covariance of all sixteen source charges}
The full physical wavefunction carrier is
$\OriginalLtwo=L^2(\mathbb R^{27};\Fermion)$, and the invariant
Hilbert subspace is $\mathcal H_{\mathrm{phys}}$ from
Chapter~\ref{ch:gauge}. If a correctly normalized Fock-vector
unitary $U_{\Fermion}(g)$ and spatial bosonic action $R_B(g)$ are
established, the combined transformation is expected to have the form
\begin{equation}\label{eq:combined-action}
 (\mathcal U_g\psi)(x)
   =U_{\Fermion}(g)\psi(R_B(g)^{-1}x),
\end{equation}
where conventions must be checked against the exact source
\texttt{bosonPullback} and \texttt{fiberAction}. An operator-valued
Hilbert--Schmidt action, even when unitary as an abstract fiber, is
not automatically the action \eqref{eq:combined-action} on the
original physical Fock-vector wavefunctions.

Let $T$ denote the actual closed sixteen-charge column on
$\mathcal H_{\mathrm{phys}}$ from Chapter~\ref{ch:hamiltonian}.
A source-faithful all-spinor covariance would be an identity
\begin{equation}\label{eq:full-charge-covariance-target}
 T\mathcal U_g=\mathcal S_g T
\end{equation}
on the literal source domain, where $\mathcal S_g$ is the corresponding
charge-output Hilbert-space unitary. Such a statement is stronger
than the separate source kinetic covariance or individual cubic
commutator lemmas: it must cover the exact original sixteen charges,
the original $\gamma$ and $\theta$ normalizations, full color sums,
Fock action, gauge restriction and original domain.

\begin{proposition}[Closed-form transport from a verified charge intertwiner;
conditional functional analysis]
\label{prop:closed-form-transport}
Suppose $\mathcal U_g$ and $\mathcal S_g$ are unitary actions on the
original physical input and original charge-output Hilbert spaces.
Suppose the original invariant smooth core is preserved and
\eqref{eq:full-charge-covariance-target} holds there. If that core
is a graph core for the closed $T$, then
$\mathcal U_g\operatorname{Dom}(T)=\operatorname{Dom}(T)$ and
\begin{equation}\label{eq:form-invariance-target}
 q(\mathcal U_gu,\mathcal U_gv)=q(u,v),\qquad
 q(u,v)=\tfrac1{16}\langle Tu,Tv\rangle.
\end{equation}
Consequently the associated self-adjoint Hamiltonian is invariant:
$\mathcal U_g\operatorname{Dom}(H)=\operatorname{Dom}(H)$ and
$H\mathcal U_g=\mathcal U_gH$ on its natural domain.
\end{proposition}
\begin{proof}[Conditional proof]
Take $u\in\operatorname{Dom}(T)$ and choose core vectors $u_n$ with
$(u_n,Tu_n)\to(u,Tu)$ in the original graph norm. Applying the unitary
pair $(\mathcal U_g,\mathcal S_g)$ preserves convergence. The core
intertwining identity puts
$(\mathcal U_gu_n,\mathcal S_gTu_n)$ in the graph of $T$. Closedness
of $T$ therefore yields
$(\mathcal U_gu,\mathcal S_gTu)$ in the graph. The same argument
with $g^{-1}$ proves equality of the transformed domain.
Unitarity of $\mathcal S_g$ yields \eqref{eq:form-invariance-target}.
The associated-vector characterization of $H$ from
Chapter~\ref{ch:hamiltonian} then transports its operator domain and
action under the invariant form. No guessed differential expression
is substituted for the original $T^*T$.
\end{proof}
\editornote{This is a conventional implication under explicit hypotheses.
The source-specific $\mathcal U_g$, physical gauge invariance, and full
sixteen-charge intertwining have not been jointly accepted; the proposition
must not be cited as a completed physical Spin(9) Hamiltonian symmetry.}

\section{From a physical symmetry to reducing isotypic sectors}
Only after a strongly continuous compact-group action on the \emph{same}
physical Hilbert space is established do its character projections
become available. For an irreducible Spin type $\lambda$, one expects
\begin{equation}\label{eq:isotype-projector}
 P_\lambda=d_\lambda\int_{G_9}
   \overline{\chi_\lambda(g)}\mathcal U_g\,dg,
\end{equation}
with normalized Haar measure. The projection must preserve the
physical closed form domain, not just the Hilbert space, if it is to
reduce the associated Hamiltonian. A typical path proves the
projected smooth invariant core is graph/form dense, then uses form
invariance and the associated-operator theorem to obtain a reducing
sector. That route must be checked against the original physical
core, the precise group topology, and the original output action.

In particular, existence of an algebraic character projector on an
independently chosen module does not prove compact resolvent in the
original physical sector. The conditional high-weight oscillator
argument in Chapter~\ref{ch:spectral} also needs the actual
$\operatorname{Spin}(8)$ branching and the correct bound for the
chosen fermionic carrier. Without those source bridges, the
subsequent radial coercivity and positive-eigenpair implications remain
conditional.

\section{Why the history of failed large proofs belongs here}
The FCP S3L record preserves accepted kinetic covariance, the
original cubic response components, and a sequence of failed large
final specializations. The failures were source elaboration and kernel
complexity phenomena. Later C48 work showed a more productive pattern:
prove source-label commutation, independently identify the original
S1J matrix, and then transport to active gamma operators in separate
strict gates. Neither development licenses a leap from finite-label
action to closed physical Hamiltonian covariance, but their contrast
shows how the dependency structure can be made tractable without
weakening the problem.

\section{Reader's checkpoint}
The original label-level Spin(9) and SU(2) actions commute, and the
result has been connected to the \emph{actual} S1J source and the
active Fock gamma action by independently qualified gates. The next
substantive scientific obstacles are stronger: fix the genuine gauge
unitaries, construct the correct physical-vector action, prove the
sixteen-charge intertwiner and obtain the consequent closed-form
and reducing-sector theorems. Those are the precise joins needed
before the compactness and oscillator machinery can address the
spectrum of the native self-adjoint Hamiltonian.
